% P53 | 6 Comparison of methods | Result: the American put on one footing
% guide: dev_notes/paragraph_guides/P53_result_the_american_put_on_one_footing.txt
Table~\ref{tbl:deskput} shows that the tree with $M=1600$ steps is within \DeskPutTreeErr\ of the extrapolated reference and takes \DeskPutTreeMs\ ms.
The list-based class needs \DeskPutClassMs\ ms for $800$ steps against \DeskPutVecMs\ ms for the vectorised pricer, and the reference itself costs \DeskPutRefMs\ ms.
The finite-difference solver on a $801\times200$ grid has an error of \DeskPdeCoarseErr\ in \DeskPdeCoarseMs\ ms, and on a $3201\times800$ grid an error of \DeskPdeFineErr\ in \DeskPdeFineMs\ ms.
In one dimension the PDE solver is therefore at least as accurate as the tree for comparable time, and it gives the boundary and the Greeks from the same grid.
Least-squares Monte Carlo has an error of \DeskLsmErr\ against a standard error of \DeskLsmSE\ and takes \DeskLsmMs\ ms, so it is the slowest method here, and its error is of the size of its own statistical noise.
The Barone--Adesi--Whaley approximation is off by \DeskBawErr\ but takes \DeskBawMs\ ms, which is why it remains useful for screening.
Least-squares Monte Carlo is worth having nonetheless, because its cost does not grow exponentially with the number of underlyings, which is the case where the tree and the grid fail.
